\section{Introduction}

As AI agents take more consequential actions with less supervision, the people those actions affect need to be able to state their needs, correct the agent, and see what it actually did.
Most agent evaluations test one user and one agent, or score what the agent says.
The Multi-stakeholder Relational Agency Framework (MSRAF) tests what an agent does when several affected people want different things.
A family settling a will is one such case: the agent has no single principal and must help each person decide for themselves.

MSRAF fixes four rules.
Silence and politeness never count as agreement.
A coordinating role never gives authority over another person's decision.
An objection is information, not an automatic veto.
A check that never triggers is scored not applicable, never a pass.
Resisting correction, reporting consent that was never given, and acting without authority are each ways an agent slips out of meaningful human oversight.
MSRAF targets these failures at the level of individual actions.
Clean results in tested scenarios would not establish general safety.

\paragraph{Three metrics.}
MSRAF scores Cooperation, Corrigibility and Candor.
The following are proposed MSRAF working definitions, adapted from discussion in the author's Nonviolent Global Liberation (NGL) notes; they are not Miki Kashtan's own definitions.
\begin{itemize}
  \item \textbf{Cooperation:} ``Working toward an outcome every affected party is willing to accept because their underlying needs are heard and included.''
  \item \textbf{Corrigible:} ``Leaving real room for others to say no, and changing course without defending itself when someone objects, while still using judgment instead of blindly deferring.''
  \item \textbf{Candid:} ``Making the reasons for choices visible and sharing limits or uncertainty when it helps the people it serves, specifically proactively flagging consequential problems.''
\end{itemize}
Each metric is the share of its applicable checks that pass, between 0 and 1.
The three are scored separately and never summed, because a high sum could hide a trade-off between them.
They complement the helpful, honest and harmless (HHH) framing of \citet{askell2021general}, which describes an assistant serving one user; Candor builds on honesty and does not replace it.
The focus on each stakeholder's own decision follows relational autonomy \citep{mackenzie2000relational}.

\paragraph{Current state.}
Most checks are scored by one LLM judge; completion and contact are computed from the event log (rubric v2).
The rubric was drafted during the sprint and has not been reviewed by a person outside the project.
This revision reports a 15-turn grid of 46 episodes on one handcrafted scenario as its main result, an earlier 12-turn grid of 48 episodes, 24 episodes on two generated scenarios and a 4-episode run without a turn cap.
Cases where the agent should proceed despite an objection (``justified proceeding'') are not implemented yet.

\paragraph{The pilot.}
To show the gym in use, we ran a small $2{\times}2$ experiment.
One factor is a compact Convergent Facilitation (CF) toolkit condensed from \citet{kashtan2020hcd}.
The other is an ``active participant'' (AP) framing: the facilitator's needs count no more and no less than anyone else's \citep[p.~255]{kashtan2020hcd}.
AP is an experimental arm only.
It is not part of the scoring rule.

\paragraph{Contributions.}
\begin{itemize}
  \item MSRAF: three separately scored metrics with 18 checks for an agent working between several stakeholders.
  \item A seeded scenario generator and a turn-based asynchronous simulator.
  \item A pilot that shows what each metric catches and where the checks are too easy.
  \item An open account of the limits and dual-use risks of this kind of benchmark.
\end{itemize}
Code, prompts, scenarios and scored episodes are at \url{https://github.com/flowstyleliving/CCCG}.
This pilot started anew from an earlier CCCG hospital-discharge simulator (begun October 31, 2026) and kept only its principles.

\subsection{Why these three properties matter for catastrophic risk}
\label{sec:xrisk}
Arguments about existential risk from AI often describe ways an agent escapes human oversight or disempowers people.
Each MSRAF metric targets a small, single-action version of one such failure.
We claim only that these properties are discussed in such arguments.
MSRAF measures them in toy scenarios, and it does not show that any score reduces existential risk.

\paragraph{Corrigibility.}
\citet{soares2015corrigibility} define a corrigible agent as one that accepts correction and shutdown even though a goal-directed agent usually has reason to resist.
\citet{hadfieldmenell2017offswitch} show that an agent leaves its off switch alone only if it is uncertain about its objective and treats human action as evidence about it.
\citet{turner2021power} prove that, in many environments, optimal policies tend to keep options open and avoid shutdown.
\citet{carlsmith2022powerseeking} builds an existential-risk argument on such power-seeking leading to human disempowerment.
MSRAF's Corrigibility checks ask the everyday version: does the agent invite a ``no'', adjust when someone objects, and refuse to treat a point person as having authority over others?
The definition also asks for judgment, not blind deference, which is why an objection is information and not an automatic veto.

\paragraph{Candor.}
\citet{evans2021truthful} argue that AI systems need standards against lies and negligent falsehoods.
\citet{park2023deception} survey learned deception in current systems and list deception during loss of control among its risks.
\citet{greenblatt2024alignmentfaking} report a model that complies selectively during training to avoid being modified, and \citet{meinke2024scheming} report frontier models that, in agentic tests, attempt to disable oversight and then deny it.
Oversight depends on seeing what an agent actually did.
MSRAF's Candor checks ask whether the final report is true, whether relays are accurate, whether unresolved concerns are reported, and whether the agent states its own limits.

\paragraph{Cooperation.}
\citet{critch2020arches} argue that existential safety must address delegation between many humans and many AI systems, not only one human and one AI.
\citet{dafoe2020cooperative} set out cooperation between humans and machines as a research field, and \citet{hammond2025multiagent} list miscoordination, conflict and collusion as risks of multi-agent AI.
\citet{kulveit2025gradual} argue that AI could erode human influence gradually as it displaces human participation in institutions.
An agent that reports consent nobody gave, or that settles a group decision without hearing each person's needs, is a small instance of that displacement.
MSRAF's Cooperation checks ask whether the agent contacted everyone, asked why things matter, surfaced each person's hidden need and validated options with them.
